\addtocontents{toc}{\protect\setcounter{tocdepth}{0}}
\chapter*{LỜI MỞ ĐẦU}\label{loimodau_intro}
\addcontentsline{toc}{chapter}{LỜI MỞ ĐẦU}
\vspace{-40pt}

\renewcommand{\thesection}{\arabic{section}} % Loại bỏ số chương (0) phía trước, chỉ hiện số thứ tự section
% ------------------------------------------
\section{Lý do chọn đề tài}
Nội dung đoạn 1
%- gõ enter 2 lần để xuống đoạn

Nội dung đoạn 2 nếu có
\section{Mục tiêu nghiên cứu}
Nội dung đoạn 1


Nội dung đoạn 2 nếu có
\section{Đối tượng và phạm vi nghiên cứu}
\begin{itemize}
    \item \textbf{Đối tượng nghiên cứu}
    
    Nội dung đoạn 1


    Nội dung đoạn 2 nếu có
    \item \textbf{Phạm vi nghiên cứu}
    
    Nội dung đoạn 1


    Nội dung đoạn 2 nếu có
\end{itemize}


\section{Phương pháp nghiên cứu}
\begin{itemize}
    \item \textbf{Phương pháp nghiên cứu lý thuyết:} Tìm hiểu các tài liệu về công nghệ [Tên công nghệ bạn dùng], các thuật toán liên quan để xây dựng cơ sở lý luận.
    \item \textbf{Phương pháp thực nghiệm:} Xây dựng ứng dụng thử nghiệm trên nền tảng [Android/Web/AI...]. Tiến hành chạy thử nghiệm với bộ dữ liệu [Tên bộ dữ liệu] để kiểm chứng.
    \item \textbf{Phương pháp phân tích và đánh giá:} So sánh hiệu năng của hệ thống mới so với các giải pháp hiện có về mặt thời gian phản hồi và độ chính xác.
\end{itemize}


\section{Cấu trúc đồ án tốt nghiệp}
\begin{itemize}
    \item \textbf{Chương 1: Tổng quan lý thuyết và cơ sở nghiên cứu} \\
    Nghiên cứu tổng quan về các công nghệ [Tên công nghệ], các giao thức và nền tảng kiến thức nền tảng làm cơ sở để thực hiện đề tài.
    
    \item \textbf{Chương 2: Phương pháp nghiên cứu} \\
    Trình bày kết quả khảo sát hiện trạng, phân tích yêu cầu (chức năng và phi chức năng), thiết kế kiến trúc hệ thống, sơ đồ logic và thiết kế cơ sở dữ liệu.
    
    \item \textbf{Chương 3: Kết quả và thảo luận} \\
   Trình bày các giao diện chức năng chính của sản phẩm.
    
    \item \textbf{Chương 4: Kết luận và đề xuất} \\
   Kết luận và đề xuất ....
     \item \textbf{Tài liệu tham khảo} 
       \item \textbf{Phụ lục (nếu có)}
\end{itemize}